\documentclass[10pt,oneside]{report}
\usepackage{soc2style}
\begin{document}
\doctitle{Trust Services Criteria Technical Reference}{AICPA criterion catalog with non-normative engineering notes}{SOC2-01}
\section*{Operating convention and boundaries}
An AICPA criterion identifies an outcome against which controls may be evaluated. The selected point-of-focus headings identify relevant characteristics; they are not discrete certification checklist items. Every owner assignment, technological mechanism and audit test in this suite is a proposed example to be tailored to the system description, commitments, risk profile, period, and service organization responsibilities.

\begin{itemize}
\item Always establish system boundaries, the in-scope trust services categories, service commitments and system requirements first.
\item For a security-only engagement, use the common criteria. For additional trust services categories, include the corresponding additional criteria, subject to properly justified applicability.
\item Document complementary user entity controls and subservice organizations; do not treat vendor assurance reports as replacements for management ownership.
\item Privacy points of focus marked [P] or [C] describe typical processor or controller relevance; assess actual responsibilities rather than treating the annotations as absolute.
\end{itemize}

\chapter{Purpose, structure, and scope}
The 2017 Trust Services Criteria (TSC) support evaluation of the design and operating effectiveness of controls over security, availability, processing integrity, confidentiality and privacy. The criteria are aligned to COSO principles and supplemented with logical and physical access, system operations, change management and risk mitigation criteria. The specific form of safeguards is a management choice responsive to risk; the criteria themselves are not a prescriptive tool inventory.

The 2022 edits revise points of focus, while the criteria stay unchanged. The source expressly permits tailored points of focus and does not require a separate assessment of every point of focus. Accordingly, the extracted focus headings below are deliberately selected illustrative excerpts rather than a complete reproduction of the many point-of-focus statements.

\subsection{Trust services categories}
\begin{itemize}
\item Security: the common criteria (CC1--CC9) are suitable without a separate S-series.
\item Availability: common criteria plus the A1.1--A1.3 series.
\item Processing Integrity: common criteria plus PI1.1--PI1.5, concerning complete, valid, accurate, timely, authorized processing.
\item Confidentiality: common criteria plus C1.1--C1.2, covering information designated confidential.
\item Privacy: common criteria plus P1.1--P8.1, addressing personal information throughout the privacy lifecycle.
\end{itemize}

\chapter{Criterion catalog with annotated engineering interpretation}
\section{Control Environment}
\Needspace{5\baselineskip}
\subsection*{CC1.1 -- Ethics and conduct}
\begin{aicpabox}
\textbf{CC1.1}. COSO Principle 1: The entity demonstrates a commitment to integrity and ethical values.\par\sourcepage{14}
\end{aicpabox}
Selected AICPA point-of-focus headings (illustrative, not exhaustive): Sets the Tone at the Top; Establishes Standards of Conduct; Evaluates Adherence to Standards of Conduct.

\begin{engineeringbox}
Publish and annually attest a code of conduct; operate an independent reporting and escalation path.
\end{engineeringbox}
Suggested accountable function: Legal / HR. Possible evidence: Approved conduct policy; attestations; case register with redactions.

Illustrative test: Select hires and incumbent population; inspect communication, acknowledgment and deviations.

\Needspace{5\baselineskip}
\subsection*{CC1.2 -- Board oversight}
\begin{aicpabox}
\textbf{CC1.2}. COSO Principle 2: The board of directors demonstrates independence from management and exercises oversight of the development and performance of internal control.\par\sourcepage{15}
\end{aicpabox}
Selected AICPA point-of-focus headings (illustrative, not exhaustive): Establishes Oversight Responsibilities; Applies Relevant Expertise; Operates Independently.

\begin{engineeringbox}
Schedule risk and control oversight sessions with documented challenges, decisions and assigned actions.
\end{engineeringbox}
Suggested accountable function: Board secretary. Possible evidence: Board or committee minutes; risk packs; action tracker.

Illustrative test: Trace material control exceptions to board reporting and closure.

\Needspace{5\baselineskip}
\subsection*{CC1.3 -- Structures and responsibility}
\begin{aicpabox}
\textbf{CC1.3}. COSO Principle 3: Management establishes, with board oversight, structures, reporting lines, and appropriate authorities and responsibilities in the pursuit of objectives.\par\sourcepage{16}
\end{aicpabox}
Selected AICPA point-of-focus headings (illustrative, not exhaustive): Considers All Structures of the Entity; Establishes Reporting Lines; Defines, Assigns, and Limits Authorities and Responsibilities.

\begin{engineeringbox}
Maintain organization chart, RACI, control ownership and vendor oversight responsibility.
\end{engineeringbox}
Suggested accountable function: GRC / Executive. Possible evidence: Org chart; role descriptions; RACI; delegations.

Illustrative test: Sample high-risk controls for assigned accountable owner and segregation.

\Needspace{5\baselineskip}
\subsection*{CC1.4 -- Competence and training}
\begin{aicpabox}
\textbf{CC1.4}. COSO Principle 4: The entity demonstrates a commitment to attract, develop, and retain competent individuals in alignment with objectives.\par\sourcepage{17}
\end{aicpabox}
Selected AICPA point-of-focus headings (illustrative, not exhaustive): Establishes Policies and Practices; Evaluates Competence and Addresses Shortcomings; Attracts, Develops, and Retains Individuals.

\begin{engineeringbox}
Use role profiles, background screening where lawful, onboarding, continuing training and succession plans.
\end{engineeringbox}
Suggested accountable function: HR / Security. Possible evidence: Role matrix; training records; screening attestations.

Illustrative test: Sample privileged roles for qualification and current training.

\Needspace{5\baselineskip}
\subsection*{CC1.5 -- Accountability and performance}
\begin{aicpabox}
\textbf{CC1.5}. COSO Principle 5: The entity holds individuals accountable for their internal control responsibilities in the pursuit of objectives.\par\sourcepage{18}
\end{aicpabox}
Selected AICPA point-of-focus headings (illustrative, not exhaustive): Establishes Performance Measures, Incentives, and Rewards; Considers Excessive Pressures; Evaluates Performance and Rewards or Disciplines Individuals.

\begin{engineeringbox}
Tie control responsibilities to performance review, corrective action and incentives with safe escalation.
\end{engineeringbox}
Suggested accountable function: HR / Compliance. Possible evidence: Review rubric; acknowledgments; corrective-action logs.

Illustrative test: Trace missed control duties to accountability and follow-up.

\section{Information and Communication}
\Needspace{5\baselineskip}
\subsection*{CC2.1 -- Quality information and inventory}
\begin{aicpabox}
\textbf{CC2.1}. COSO Principle 13: The entity obtains or generates and uses relevant, quality information to support the functioning of internal control.\par\sourcepage{19}
\end{aicpabox}
Selected AICPA point-of-focus headings (illustrative, not exhaustive): Identifies Information Requirements; Captures Internal and External Sources of Data; Processes Relevant Data Into Information.

\begin{engineeringbox}
Inventory services, repositories, assets, vendors and data flows with ownership, classification and periodic reconciliation.
\end{engineeringbox}
Suggested accountable function: Asset / Data Governance. Possible evidence: CMDB; data-flow diagram; classification register; reconciliation logs.

Illustrative test: Reconcile sampled production resources against inventory and classifications.

\Needspace{5\baselineskip}
\subsection*{CC2.2 -- Internal communications}
\begin{aicpabox}
\textbf{CC2.2}. COSO Principle 14: The entity internally communicates information, including objectives and responsibilities for internal control, necessary to support the functioning of internal control.\par\sourcepage{21}
\end{aicpabox}
Selected AICPA point-of-focus headings (illustrative, not exhaustive): Communicates Internal Control Information; Communicates With the Board of Directors; Provides Separate Communication Lines.

\begin{engineeringbox}
Publish control objectives, incident reporting channels, policy changes and security awareness to relevant personnel.
\end{engineeringbox}
Suggested accountable function: GRC / Comms. Possible evidence: Communications log; trainings; incident channel instructions.

Illustrative test: Sample changes and staff for timely notification and comprehension.

\Needspace{5\baselineskip}
\subsection*{CC2.3 -- External communications}
\begin{aicpabox}
\textbf{CC2.3}. COSO Principle 15: The entity communicates with external parties regarding matters affecting the functioning of internal control.\par\sourcepage{22}
\end{aicpabox}
Selected AICPA point-of-focus headings (illustrative, not exhaustive): Communicates to External Parties; Enables Inbound Communications; Communicates With the Board of Directors.

\begin{engineeringbox}
Define accurate customer, regulator, vendor and incident contact channels and commitments.
\end{engineeringbox}
Suggested accountable function: Legal / Customer Trust. Possible evidence: Service agreements; trust communications; contact playbook.

Illustrative test: Select applicable customer commitments and inspect external communication.

\section{Risk Assessment}
\Needspace{5\baselineskip}
\subsection*{CC3.1 -- Objectives and risk appetite}
\begin{aicpabox}
\textbf{CC3.1}. COSO Principle 6: The entity specifies objectives with sufficient clarity to enable the identification and assessment of risks relating to objectives.\par\sourcepage{24}
\end{aicpabox}
Selected AICPA point-of-focus headings (illustrative, not exhaustive): Reflects Management's Choices; Considers Tolerances for Risk; Includes Operations and Financial Performance Goals.

\begin{engineeringbox}
Record system objectives, commitments, risk tolerances and system boundaries before selecting safeguards.
\end{engineeringbox}
Suggested accountable function: Executive / Risk. Possible evidence: Risk appetite; system description; SLA register.

Illustrative test: Trace each material commitment into risk identification and controls.

\Needspace{5\baselineskip}
\subsection*{CC3.2 -- Risk identification and analysis}
\begin{aicpabox}
\textbf{CC3.2}. COSO Principle 7: The entity identifies risks to the achievement of its objectives across the entity and analyzes risks as a basis for determining how the risks should be managed.\par\sourcepage{27}
\end{aicpabox}
Selected AICPA point-of-focus headings (illustrative, not exhaustive): Includes Entity, Subsidiary, Division, Operating Unit, and Functional Levels; Analyzes Internal and External Factors; Involves Appropriate Levels of Management.

\begin{engineeringbox}
Maintain risk inventory, threat scenarios, likelihood and impact, treatment and periodic re-evaluation.
\end{engineeringbox}
Suggested accountable function: Enterprise Risk. Possible evidence: Risk register; workshop results; treatment approvals.

Illustrative test: Sample assets/services and verify risk analysis and owners.

\Needspace{5\baselineskip}
\subsection*{CC3.3 -- Fraud risk}
\begin{aicpabox}
\textbf{CC3.3}. COSO Principle 8: The entity considers the potential for fraud in assessing risks to the achievement of objectives.\par\sourcepage{28}
\end{aicpabox}
Selected AICPA point-of-focus headings (illustrative, not exhaustive): Considers Various Types of Fraud; Assesses Incentives and Pressures; Assesses Opportunities.

\begin{engineeringbox}
Model insider misuse, collusion, credential abuse and financial misrepresentation; design anti-fraud controls.
\end{engineeringbox}
Suggested accountable function: Risk / Internal Audit. Possible evidence: Fraud scenarios; SoD review; whistleblower handling.

Illustrative test: Trace identified fraud scenarios to mitigations and monitoring.

\Needspace{5\baselineskip}
\subsection*{CC3.4 -- Significant change assessment}
\begin{aicpabox}
\textbf{CC3.4}. COSO Principle 9: The entity identifies and assesses changes that could significantly impact the system of internal control.\par\sourcepage{29}
\end{aicpabox}
Selected AICPA point-of-focus headings (illustrative, not exhaustive): Assesses Changes in the External Environment; Assesses Changes in the Business Model; Assesses Changes in Leadership.

\begin{engineeringbox}
Trigger risk reviews for material architecture, acquisition, regulation and vendor changes.
\end{engineeringbox}
Suggested accountable function: Architecture / Risk. Possible evidence: Change intake; impact assessments; risk decisions.

Illustrative test: Sample major changes and assess whether risk review preceded deployment.

\section{Monitoring Activities}
\Needspace{5\baselineskip}
\subsection*{CC4.1 -- Control monitoring}
\begin{aicpabox}
\textbf{CC4.1}. COSO Principle 16: The entity selects, develops, and performs ongoing and/or separate evaluations to ascertain whether the components of internal control are present and functioning.\par\sourcepage{30}
\end{aicpabox}
Selected AICPA point-of-focus headings (illustrative, not exhaustive): Considers a Mix of Ongoing and Separate Evaluations; Considers Rate of Change; Establishes Baseline Understanding.

\begin{engineeringbox}
Maintain a monitoring plan mixing continuous telemetry and independent targeted reviews.
\end{engineeringbox}
Suggested accountable function: Control Assurance. Possible evidence: Control monitoring schedule; dashboards; findings.

Illustrative test: Compare monitoring coverage and executed reviews to control inventory.

\Needspace{5\baselineskip}
\subsection*{CC4.2 -- Deficiency communication}
\begin{aicpabox}
\textbf{CC4.2}. COSO Principle 17: The entity evaluates and communicates internal control deficiencies in a timely manner to those parties responsible for taking corrective action, including senior management and the board of directors, as appropriate.\par\sourcepage{31}
\end{aicpabox}
Selected AICPA point-of-focus headings (illustrative, not exhaustive): Assesses Results; Communicates Deficiencies; Monitors Corrective Action.

\begin{engineeringbox}
Grade deficiencies, route them to accountable owners, track deadlines and validate closure.
\end{engineeringbox}
Suggested accountable function: GRC / Risk. Possible evidence: Findings log; escalation tickets; validation artifacts.

Illustrative test: Sample overdue issues for escalation and independent closure verification.

\section{Control Activities}
\Needspace{5\baselineskip}
\subsection*{CC5.1 -- Risk-based control activities}
\begin{aicpabox}
\textbf{CC5.1}. COSO Principle 10: The entity selects and develops control activities that contribute to the mitigation of risks to the achievement of objectives to acceptable levels.\par\sourcepage{31}
\end{aicpabox}
Selected AICPA point-of-focus headings (illustrative, not exhaustive): Integrates With Risk Assessment; Considers Entity-Specific Factors; Determines Relevant Business Processes.

\begin{engineeringbox}
Choose preventive/detective and manual/automated controls tied to documented risks.
\end{engineeringbox}
Suggested accountable function: Risk / Control Owners. Possible evidence: Risk-control matrix; approvals; control narratives.

Illustrative test: Walk through sampled risks and inspect implemented control activity.

\Needspace{5\baselineskip}
\subsection*{CC5.2 -- Technology general controls}
\begin{aicpabox}
\textbf{CC5.2}. COSO Principle 11: The entity also selects and develops general control activities over technology to support the achievement of objectives.\par\sourcepage{32}
\end{aicpabox}
Selected AICPA point-of-focus headings (illustrative, not exhaustive): Establishes Relevant Technology Infrastructure Control Activities; Establishes Relevant Security Management Process Controls Activities.

\begin{engineeringbox}
Establish technology operations, infrastructure, access, development and maintenance governance.
\end{engineeringbox}
Suggested accountable function: IT / Platform. Possible evidence: ITGC inventory; platform standards; administrative review.

Illustrative test: Test technology dependencies of critical automated controls.

\Needspace{5\baselineskip}
\subsection*{CC5.3 -- Policy deployment}
\begin{aicpabox}
\textbf{CC5.3}. COSO Principle 12: The entity deploys control activities through policies that establish what is expected and in procedures that put policies into action.\par\sourcepage{33}
\end{aicpabox}
Selected AICPA point-of-focus headings (illustrative, not exhaustive): Performs in a Timely Manner; Takes Corrective Action; Performs Using Competent Personnel.

\begin{engineeringbox}
Convert approved policy into assigned, timed operating procedures and exceptions governance.
\end{engineeringbox}
Suggested accountable function: GRC / Operational Owners. Possible evidence: Policies; procedures; schedules; exception register.

Illustrative test: Select policies and trace to accountable operational executions.

\section{Logical and Physical Access Controls}
\Needspace{5\baselineskip}
\subsection*{CC6.1 -- Logical access architecture}
\begin{aicpabox}
\textbf{CC6.1}. The entity implements logical access security software, infrastructure, and architectures over protected information assets to protect them from security events to meet the entity's objectives.\par\sourcepage{34}
\end{aicpabox}
Selected AICPA point-of-focus headings (illustrative, not exhaustive): Identifies and Manages the Inventory of Information Assets; Assesses New Architectures; Restricts Logical Access.

\begin{engineeringbox}
Federate identities; protect service identities; enforce MFA, RBAC and secured identity architecture.
\end{engineeringbox}
Suggested accountable function: IAM / Platform. Possible evidence: IAM configuration; architecture; identity inventory.

Illustrative test: Inspect identities and authentication design across in-scope assets.

\Needspace{5\baselineskip}
\subsection*{CC6.2 -- Provisioning and deprovisioning}
\begin{aicpabox}
\textbf{CC6.2}. Prior to issuing system credentials and granting system access, the entity registers and authorizes new internal and external users whose access is administered by the entity. For those users whose access is administered by the entity, user system credentials are removed when user access is no longer authorized.\par\sourcepage{36}
\end{aicpabox}
Selected AICPA point-of-focus headings (illustrative, not exhaustive): Reviews Validity of Access Credentials; Prevents the Use of Credentials When No Longer Valid.

\begin{engineeringbox}
Bind user creation, unique identities and timely removal to approved joiner-mover-leaver workflow.
\end{engineeringbox}
Suggested accountable function: IAM / HR. Possible evidence: Access requests; approval log; HR termination feed.

Illustrative test: Sample additions and leavers from source population to access outcomes.

\Needspace{5\baselineskip}
\subsection*{CC6.3 -- Least privilege review}
\begin{aicpabox}
\textbf{CC6.3}. The entity authorizes, modifies, or removes access to data, software, functions, and other protected information assets based on roles, responsibilities, or the system design and changes, giving consideration to the concepts of least privilege and segregation of duties, to meet the entity's objectives.\par\sourcepage{36}
\end{aicpabox}
Selected AICPA point-of-focus headings (illustrative, not exhaustive): Creates or Modifies Access to Protected Information Assets; Removes Access to Protected Information Assets; Uses Access Control Structures.

\begin{engineeringbox}
Use least-privilege roles, approval for privilege elevation, periodic access recertification and revocation.
\end{engineeringbox}
Suggested accountable function: IAM / Application Owners. Possible evidence: Access review signoffs; role policy; permissions exports.

Illustrative test: Trace sampled privileged and ordinary access to approved role requirements.

\Needspace{5\baselineskip}
\subsection*{CC6.4 -- Physical access}
\begin{aicpabox}
\textbf{CC6.4}. The entity restricts physical access to facilities and protected information assets (for example, data center facilities, backup media storage, and other sensitive locations) to authorized personnel to meet the entity's objectives.\par\sourcepage{37}
\end{aicpabox}
Selected AICPA point-of-focus headings (illustrative, not exhaustive): Creates or Modifies Physical Access; Removes Physical Access; Recovers Physical Devices.

\begin{engineeringbox}
Restrict facilities and backup-media access with appropriate badges, logs and outsourced control coverage.
\end{engineeringbox}
Suggested accountable function: Facilities / Cloud Provider. Possible evidence: Badge access review; facility records; subservice reports.

Illustrative test: Inspect applicable facility access and cloud provider responsibility boundary.

\Needspace{5\baselineskip}
\subsection*{CC6.5 -- Secure asset disposal}
\begin{aicpabox}
\textbf{CC6.5}. The entity discontinues logical and physical protections over physical assets only after the ability to read or recover data and software from those assets has been diminished and is no longer required to meet the entity's objectives.\par\sourcepage{37}
\end{aicpabox}
Selected AICPA point-of-focus headings (illustrative, not exhaustive): Removes Data and Software for Disposal.

\begin{engineeringbox}
Erase, destroy or make unrecoverable data and software before disposition; retain chain of custody.
\end{engineeringbox}
Suggested accountable function: IT Asset Management. Possible evidence: Sanitization records; disposal certificates; asset ledger.

Illustrative test: Sample decommissioned assets to verify destruction and approval.

\Needspace{5\baselineskip}
\subsection*{CC6.6 -- Perimeter and external threat defenses}
\begin{aicpabox}
\textbf{CC6.6}. The entity implements logical access security measures to protect against threats from sources outside its system boundaries.\par\sourcepage{38}
\end{aicpabox}
Selected AICPA point-of-focus headings (illustrative, not exhaustive): Restricts Access; Protects Identification and Authentication Credentials; Requires Additional Authentication or Credentials.

\begin{engineeringbox}
Protect exposed interfaces with restricted network paths, identity controls and hardening.
\end{engineeringbox}
Suggested accountable function: Network / Cloud. Possible evidence: Firewall rules; segmentation map; WAF/config reviews.

Illustrative test: Select exposed services and test approved ingress and egress posture.

\Needspace{5\baselineskip}
\subsection*{CC6.7 -- Secure data movement}
\begin{aicpabox}
\textbf{CC6.7}. The entity restricts the transmission, movement, and removal of information to authorized internal and external users and processes, and protects it during transmission, movement, or removal to meet the entity's objectives.\par\sourcepage{38}
\end{aicpabox}
Selected AICPA point-of-focus headings (illustrative, not exhaustive): Restricts the Ability to Perform Transmission; Protects Removal Media; Protects Endpoint Devices.

\begin{engineeringbox}
Classify data and restrict export, transmission and removable media; enforce suitable encryption.
\end{engineeringbox}
Suggested accountable function: Data Security. Possible evidence: DLP policy; TLS configurations; export approvals.

Illustrative test: Sample sensitive data transfers for authorized recipients and protection.

\Needspace{5\baselineskip}
\subsection*{CC6.8 -- Malicious software prevention}
\begin{aicpabox}
\textbf{CC6.8}. The entity implements controls to prevent or detect and act upon the introduction of unauthorized or malicious software to meet the entity's objectives.\par\sourcepage{39}
\end{aicpabox}
Selected AICPA point-of-focus headings (illustrative, not exhaustive): Restricts Installation and Modification of Application and Software; Detects Unauthorized Changes to Software and Configuration Parameters; Uses a Defined Change Control Process.

\begin{engineeringbox}
Prevent unapproved software and monitor malware on supported endpoints and workloads.
\end{engineeringbox}
Suggested accountable function: Endpoint / Platform. Possible evidence: EDR coverage; quarantine events; allowlist policy.

Illustrative test: Reconcile endpoint population to protection status and exceptions.

\section{System Operations}
\Needspace{5\baselineskip}
\subsection*{CC7.1 -- Vulnerability and drift detection}
\begin{aicpabox}
\textbf{CC7.1}. To meet its objectives, the entity uses detection and monitoring procedures to identify (1) changes to configurations that result in the introduction of new vulnerabilities, and (2) susceptibilities to newly discovered vulnerabilities.\par\sourcepage{40}
\end{aicpabox}
Selected AICPA point-of-focus headings (illustrative, not exhaustive): Uses Defined Configuration Standards; Monitors Infrastructure and Software; Implements Change-Detection Mechanisms.

\begin{engineeringbox}
Monitor vulnerability disclosures, baseline drift and asset exposure; prioritize and track remediation.
\end{engineeringbox}
Suggested accountable function: Vulnerability Management. Possible evidence: Scan results; configuration diffs; risk acceptance tickets.

Illustrative test: Sample vulnerable assets and follow discovery, ownership and remediation.

\Needspace{5\baselineskip}
\subsection*{CC7.2 -- Anomaly monitoring}
\begin{aicpabox}
\textbf{CC7.2}. The entity monitors system components and the operation of those components for anomalies that are indicative of malicious acts, natural disasters, and errors affecting the entity's ability to meet its objectives; anomalies are analyzed to determine whether they represent security events.\par\sourcepage{40}
\end{aicpabox}
Selected AICPA point-of-focus headings (illustrative, not exhaustive): Implements Detection Policies, Procedures, and Tools; Designs Detection Measures; Implements Filters to Analyze Anomalies.

\begin{engineeringbox}
Collect and correlate infrastructure, application and identity signals with ownership and defined review.
\end{engineeringbox}
Suggested accountable function: SOC / SRE. Possible evidence: Detection rules; SIEM alerts; on-call histories.

Illustrative test: Select detections and inspect triage and tuning evidence.

\Needspace{5\baselineskip}
\subsection*{CC7.3 -- Security event evaluation}
\begin{aicpabox}
\textbf{CC7.3}. The entity evaluates security events to determine whether they could or have resulted in a failure of the entity to meet its objectives (security incidents) and, if so, takes actions to prevent or address such failures.\par\sourcepage{41}
\end{aicpabox}
Selected AICPA point-of-focus headings (illustrative, not exhaustive): Responds to Security Incidents; Communicates and Reviews Detected Security Events; Develops and Implements Procedures to Analyze Security Incidents.

\begin{engineeringbox}
Classify events, assess impact, escalate incidents and preserve investigation records.
\end{engineeringbox}
Suggested accountable function: SOC / Incident Response. Possible evidence: Triage cases; severity rubric; incident timeline.

Illustrative test: Sample events through classification and incident decision records.

\Needspace{5\baselineskip}
\subsection*{CC7.4 -- Incident response}
\begin{aicpabox}
\textbf{CC7.4}. The entity responds to identified security incidents by executing a defined incident-response program to understand, contain, remediate, and communicate security incidents, as appropriate.\par\sourcepage{42}
\end{aicpabox}
Selected AICPA point-of-focus headings (illustrative, not exhaustive): Assigns Roles and Responsibilities; Contains and Responds to Security Incidents; Mitigates Ongoing Security Incidents.

\begin{engineeringbox}
Practice detection, containment, eradication, coordination, communications and post-incident review.
\end{engineeringbox}
Suggested accountable function: Security Operations. Possible evidence: IR plan; tabletop results; incident cases.

Illustrative test: Walk an incident through roles, approvals, communication and remediation.

\Needspace{5\baselineskip}
\subsection*{CC7.5 -- Security incident recovery}
\begin{aicpabox}
\textbf{CC7.5}. The entity identifies, develops, and implements activities to recover from identified security incidents.\par\sourcepage{44}
\end{aicpabox}
Selected AICPA point-of-focus headings (illustrative, not exhaustive): Restores the Affected Environment; Communicates Information About the Incident; Determines Root Cause of the Incident.

\begin{engineeringbox}
Define restoration sequencing, integrity checks and post-incident validation and lessons learned.
\end{engineeringbox}
Suggested accountable function: SRE / Business Continuity. Possible evidence: Recovery playbooks; restore tests; PIR action list.

Illustrative test: Trace recovered services to tested restoration and business acceptance.

\section{Change Management}
\Needspace{5\baselineskip}
\subsection*{CC8.1 -- Controlled change management}
\begin{aicpabox}
\textbf{CC8.1}. The entity authorizes, designs, develops or acquires, configures, documents, tests, approves, and implements changes to infrastructure, data, software, and procedures to meet its objectives.\par\sourcepage{45}
\end{aicpabox}
Selected AICPA point-of-focus headings (illustrative, not exhaustive): Manages Changes Throughout the System Life Cycle; Authorizes Changes; Designs and Develops Changes.

\begin{engineeringbox}
Use PR reviews, required checks, approvals, release provenance, rollback and emergency change handling.
\end{engineeringbox}
Suggested accountable function: Engineering / CAB. Possible evidence: Pull requests; checks; releases; change tickets.

Illustrative test: Sample normal and emergency production changes end-to-end.

\section{Risk Mitigation}
\Needspace{5\baselineskip}
\subsection*{CC9.1 -- Business disruption mitigation}
\begin{aicpabox}
\textbf{CC9.1}. The entity identifies, selects, and develops risk mitigation activities for risks arising from potential business disruptions.\par\sourcepage{47}
\end{aicpabox}
Selected AICPA point-of-focus headings (illustrative, not exhaustive): Considers Mitigation of Risks of Business Disruption; Considers the Use of Insurance to Mitigate Financial Impact Risks.

\begin{engineeringbox}
Define plausible disruption scenarios and proportionate mitigation including alternates and financial options.
\end{engineeringbox}
Suggested accountable function: BCM / Risk. Possible evidence: BIA; continuity response plan; mitigation decisions.

Illustrative test: Trace business-impact findings to mitigations and assigned owners.

\Needspace{5\baselineskip}
\subsection*{CC9.2 -- Third-party risk}
\begin{aicpabox}
\textbf{CC9.2}. The entity assesses and manages risks associated with vendors and business partners.\par\sourcepage{48}
\end{aicpabox}
Selected AICPA point-of-focus headings (illustrative, not exhaustive): Establishes Requirements for Vendor and Business Partner Engagements; Identifies Vulnerabilities; Assesses Vendor and Business Partner Risks.

\begin{engineeringbox}
Inventory material vendors and assess security, contracts, shared responsibilities and ongoing performance.
\end{engineeringbox}
Suggested accountable function: Third-Party Risk. Possible evidence: Vendor register; diligence; contracts; monitoring evidence.

Illustrative test: Sample critical providers for due diligence and renewal reassessment.

\section{Additional Criteria for Availability}
\Needspace{5\baselineskip}
\subsection*{A1.1 -- Capacity management}
\begin{aicpabox}
\textbf{A1.1}. The entity maintains, monitors, and evaluates current processing capacity and use of system components (infrastructure, data, and software) to manage capacity demand and to enable the implementation of additional capacity to help meet its objectives.\par\sourcepage{50}
\end{aicpabox}
Selected AICPA point-of-focus headings (illustrative, not exhaustive): Measures Current Usage; Forecasts Capacity; Makes Changes Based on Forecasts.

\begin{engineeringbox}
Track utilization and saturation, forecast load and establish thresholds for scaling and architecture changes.
\end{engineeringbox}
Suggested accountable function: SRE / Capacity. Possible evidence: Capacity dashboards; forecast; scaling decisions.

Illustrative test: Recompute capacity headroom and inspect change actions on breaches.

\Needspace{5\baselineskip}
\subsection*{A1.2 -- Environmental and backup protections}
\begin{aicpabox}
\textbf{A1.2}. The entity authorizes, designs, develops or acquires, implements, operates, approves, maintains, and monitors environmental protections, software, data backup processes, and recovery infrastructure to meet its objectives.\par\sourcepage{50}
\end{aicpabox}
Selected AICPA point-of-focus headings (illustrative, not exhaustive): Identifies Environmental Threats; Designs Detection Measures; Implements and Maintains Environmental Protection Mechanisms.

\begin{engineeringbox}
Implement backup, power/environmental protection, recovery infrastructure, restoration plans and monitoring.
\end{engineeringbox}
Suggested accountable function: Infrastructure / BCM. Possible evidence: Backup jobs; retention settings; facility/provider evidence.

Illustrative test: Sample backups for success, protected storage and recovery dependencies.

\Needspace{5\baselineskip}
\subsection*{A1.3 -- Recovery testing}
\begin{aicpabox}
\textbf{A1.3}. The entity tests recovery plan procedures supporting system recovery to meet its objectives.\par\sourcepage{52}
\end{aicpabox}
Selected AICPA point-of-focus headings (illustrative, not exhaustive): Implements Business Continuity Plan Testing; Tests Integrity and Completeness of Backup Data.

\begin{engineeringbox}
Periodically exercise restoration procedures, including integrity checks, time objectives and follow-up.
\end{engineeringbox}
Suggested accountable function: BCM / SRE. Possible evidence: DR exercise plan; restore logs; issue remediation.

Illustrative test: Reperform selected recovery steps and compare outcome with defined objectives.

\section{Additional Criteria for Confidentiality}
\Needspace{5\baselineskip}
\subsection*{C1.1 -- Confidential information life cycle}
\begin{aicpabox}
\textbf{C1.1}. The entity identifies and maintains confidential information to meet the entity's objectives related to confidentiality.\par\sourcepage{52}
\end{aicpabox}
Selected AICPA point-of-focus headings (illustrative, not exhaustive): Defines and Identifies Confidential information; Retains Confidential Information; Protects Confidential Information From Destruction.

\begin{engineeringbox}
Identify confidential datasets, restrict permitted use, retain appropriately and review protection needs.
\end{engineeringbox}
Suggested accountable function: Data Governance. Possible evidence: Classification inventory; data access matrix; retention map.

Illustrative test: Sample data classes from creation through storage and access controls.

\Needspace{5\baselineskip}
\subsection*{C1.2 -- Confidential disposal}
\begin{aicpabox}
\textbf{C1.2}. The entity disposes of confidential information to meet the entity's objectives related to confidentiality.\par\sourcepage{52}
\end{aicpabox}
Selected AICPA point-of-focus headings (illustrative, not exhaustive): Identifies Confidential Information for Destruction; Destroys Confidential Information.

\begin{engineeringbox}
Trigger approved deletion, media destruction and vendor disposition at the end of retention.
\end{engineeringbox}
Suggested accountable function: Data Governance / IT. Possible evidence: Retention schedule; deletion jobs; vendor certificates.

Illustrative test: Sample expired records and verify timely irreversible disposition.

\section{Additional Criteria for Processing Integrity}
\Needspace{5\baselineskip}
\subsection*{PI1.1 -- Processing specifications}
\begin{aicpabox}
\textbf{PI1.1}. The entity obtains or generates, uses, and communicates relevant, quality information regarding the objectives related to processing, including definitions of data processed and product and service specifications, to support the use of products and services.\par\sourcepage{53}
\end{aicpabox}
Selected AICPA point-of-focus headings (illustrative, not exhaustive): Defines Data Necessary to Support a Product or Service; Defines Information Necessary to Support the Use of a Good or Product.

\begin{engineeringbox}
Define data contracts, schemas, field semantics, functional outcomes and user-facing specifications.
\end{engineeringbox}
Suggested accountable function: Product / Data Engineering. Possible evidence: Data dictionary; schema registry; acceptance criteria.

Illustrative test: Trace specifications to implemented validations and published outputs.

\Needspace{5\baselineskip}
\subsection*{PI1.2 -- Input completeness and accuracy}
\begin{aicpabox}
\textbf{PI1.2}. The entity implements policies and procedures over system inputs, including controls over completeness and accuracy, to result in products, services, and reporting to meet the entity's objectives.\par\sourcepage{54}
\end{aicpabox}
Selected AICPA point-of-focus headings (illustrative, not exhaustive): Defines Characteristics of Processing Inputs; Evaluates Processing Inputs; Creates and Maintains Records of System Inputs.

\begin{engineeringbox}
Validate required fields, types, authorization, duplicates and completeness before processing.
\end{engineeringbox}
Suggested accountable function: Data Engineering. Possible evidence: Validation rules; rejects; ingestion logs.

Illustrative test: Inject invalid and duplicate data in controlled test and inspect handling.

\Needspace{5\baselineskip}
\subsection*{PI1.3 -- Processing correctness}
\begin{aicpabox}
\textbf{PI1.3}. The entity implements policies and procedures over system processing to result in products, services, and reporting to meet the entity's objectives.\par\sourcepage{54}
\end{aicpabox}
Selected AICPA point-of-focus headings (illustrative, not exhaustive): Defines Processing Specifications; Defines Processing Activities; Detects and Corrects Processing or Production Activity Errors.

\begin{engineeringbox}
Use idempotence, error handling, reconciliation and traceability for transformations.
\end{engineeringbox}
Suggested accountable function: Platform / Engineering. Possible evidence: Job runs; test results; reconciliation outcomes.

Illustrative test: Trace representative inputs through transformations and exception queues.

\Needspace{5\baselineskip}
\subsection*{PI1.4 -- Output distribution}
\begin{aicpabox}
\textbf{PI1.4}. The entity implements policies and procedures to make available or deliver output completely, accurately, and timely in accordance with specifications to meet the entity's objectives.\par\sourcepage{55}
\end{aicpabox}
Selected AICPA point-of-focus headings (illustrative, not exhaustive): Protects Output; Distributes Output Only to Intended Parties; Distributes Output Completely and Accurately.

\begin{engineeringbox}
Validate outputs, authorize recipients, preserve integrity and meet delivery expectations.
\end{engineeringbox}
Suggested accountable function: Product / Data Engineering. Possible evidence: Delivery manifests; approvals; integrity checks.

Illustrative test: Trace approved outputs to recipients, timing and end-to-end completeness.

\Needspace{5\baselineskip}
\subsection*{PI1.5 -- Processing records retention}
\begin{aicpabox}
\textbf{PI1.5}. The entity implements policies and procedures to store inputs, items in processing, and outputs completely, accurately, and timely in accordance with system specifications to meet the entity's objectives.\par\sourcepage{56}
\end{aicpabox}
Selected AICPA point-of-focus headings (illustrative, not exhaustive): Protects Stored Items; Archives and Protects System Records; Stores Data Completely and Accurately.

\begin{engineeringbox}
Retain and protect input, intermediate and output records per system specification.
\end{engineeringbox}
Suggested accountable function: Data Engineering / Records. Possible evidence: Retention policies; transaction logs; storage tests.

Illustrative test: Sample lifecycle records for accuracy, recoverability and defined retention.

\section{Additional Criteria for Privacy: Notice}
\Needspace{5\baselineskip}
\subsection*{P1.1 -- Privacy notice}
\begin{aicpabox}
\textbf{P1.1}. The entity provides notice to data subjects about its privacy practices to meet the entity's objectives related to privacy. The notice is updated and communicated to data subjects in a timely manner for changes to the entity's privacy practices, including changes in the use of personal information, to meet the entity's objectives related to privacy.\par\sourcepage{57}
\end{aicpabox}
Selected AICPA point-of-focus headings (illustrative, not exhaustive): Communicates to Data Subjects [C]; Provides Notice to Data Subjects [C]; Covers Entities and Activities in Notice [C].

\begin{engineeringbox}
Maintain accurate notices mapped to actual data processing and communicate material updates.
\end{engineeringbox}
Suggested accountable function: Privacy / Legal. Possible evidence: Privacy notice versions; publication history; data map.

Illustrative test: Compare sampled collection paths with notice and update history.

\section{Additional Criteria for Privacy: Choice and Consent}
\Needspace{5\baselineskip}
\subsection*{P2.1 -- Choice and consent}
\begin{aicpabox}
\textbf{P2.1}. The entity communicates choices available regarding the collection, use, retention, disclosure, and disposal of personal information to the data subjects and the consequences, if any, of each choice. Explicit consent for the collection, use, retention, disclosure, and disposal of personal information is obtained from data subjects or other authorized persons, if required. Such consent is obtained only for the intended purpose of the information to meet the entity's objectives related to privacy. The entity's basis for determining implicit consent for the collection, use, retention, disclosure, and disposal of personal information is documented.\par\sourcepage{58}
\end{aicpabox}
Selected AICPA point-of-focus headings (illustrative, not exhaustive): Communicates to Data Subjects [C]; Communicates Consequences of Denying or Withdrawing Consent [C]; Obtains Implicit or Explicit Consent [C].

\begin{engineeringbox}
Capture consent or lawful choice decisions with purpose, version, timestamp and withdrawal workflow.
\end{engineeringbox}
Suggested accountable function: Privacy / Product. Possible evidence: Consent ledger; UX screens; withdrawal logs.

Illustrative test: Sample preferences to downstream enforcement and change history.

\section{Additional Criteria for Privacy: Collection}
\Needspace{5\baselineskip}
\subsection*{P3.1 -- Purpose-limited collection}
\begin{aicpabox}
\textbf{P3.1}. Personal information is collected consistent with the entity's objectives related to privacy.\par\sourcepage{59}
\end{aicpabox}
Selected AICPA point-of-focus headings (illustrative, not exhaustive): Limits the Collection of Personal Information [P][C]; Collects Information by Fair and Lawful Means [P][C]; Collects Information From Reliable Sources [P][C].

\begin{engineeringbox}
Map personal information fields to authorized purposes and minimize collection by default.
\end{engineeringbox}
Suggested accountable function: Privacy / Product. Possible evidence: Data inventory; form design; API schemas.

Illustrative test: Compare collected fields with documented purpose and approved scope.

\Needspace{5\baselineskip}
\subsection*{P3.2 -- Explicit consent collection}
\begin{aicpabox}
\textbf{P3.2}. For information requiring explicit consent, the entity communicates the need for such consent as well as the consequences of a failure to provide consent for the request for personal information and obtains the consent prior to the collection of the information to meet the entity's objectives related to privacy.\par\sourcepage{60}
\end{aicpabox}
Selected AICPA point-of-focus headings (illustrative, not exhaustive): Informs Data Subjects of Consequences of Failure to Provide Consent [C]; Documents Explicit Consent to Retain Information [C].

\begin{engineeringbox}
Where explicit consent is required, communicate consequences and preserve affirmative evidence.
\end{engineeringbox}
Suggested accountable function: Privacy / Legal. Possible evidence: Consent artifacts; notices; audit logs.

Illustrative test: Sample transactions requiring explicit consent for recorded decision.

\section{Additional Criteria for Privacy: Use, Retention, and Disposal}
\Needspace{5\baselineskip}
\subsection*{P4.1 -- Purpose limitation}
\begin{aicpabox}
\textbf{P4.1}. The entity limits the use of personal information to the purposes identified in the entity's objectives related to privacy.\par\sourcepage{60}
\end{aicpabox}
Selected AICPA point-of-focus headings (illustrative, not exhaustive): Uses Personal Information for Intended Purposes [P][C].

\begin{engineeringbox}
Implement use restrictions, lineage and secondary-use approvals for personal information.
\end{engineeringbox}
Suggested accountable function: Privacy / Data Owners. Possible evidence: Purpose map; access policy; data-use approvals.

Illustrative test: Trace secondary uses to approved purpose and technical enforcement.

\Needspace{5\baselineskip}
\subsection*{P4.2 -- Personal-information retention}
\begin{aicpabox}
\textbf{P4.2}. The entity retains personal information consistent with the entity's objectives related to privacy.\par\sourcepage{60}
\end{aicpabox}
Selected AICPA point-of-focus headings (illustrative, not exhaustive): Retains Personal Information [P][C]; Protects Personal Information [P][C].

\begin{engineeringbox}
Define retention clocks, lawful holds, processing-specific schedules and automated review.
\end{engineeringbox}
Suggested accountable function: Privacy / Records. Possible evidence: Retention schedule; legal holds; purge reports.

Illustrative test: Sample aged records against approved retention and legal holds.

\Needspace{5\baselineskip}
\subsection*{P4.3 -- Personal-information disposal}
\begin{aicpabox}
\textbf{P4.3}. The entity securely disposes of personal information to meet the entity's objectives related to privacy.\par\sourcepage{61}
\end{aicpabox}
Selected AICPA point-of-focus headings (illustrative, not exhaustive): Captures, Identifies, and Flags Requests for Deletion [P][C]; Disposes of, Destroys, and Redacts Personal Information [P][C]; Destroys Personal Information [P][C].

\begin{engineeringbox}
Implement verifiable deletion across primary stores, replicas and processor handoffs.
\end{engineeringbox}
Suggested accountable function: Privacy / Data Engineering. Possible evidence: Erasure tickets; deletion evidence; vendor confirmations.

Illustrative test: Trace selected deletion requests through complete affected systems.

\section{Additional Criteria for Privacy: Access}
\Needspace{5\baselineskip}
\subsection*{P5.1 -- Data subject access}
\begin{aicpabox}
\textbf{P5.1}. The entity grants identified and authenticated data subjects the ability to access their stored personal information for review and, upon request, provides physical or electronic copies of that information to data subjects to meet the entity's objectives related to privacy. If access is denied, data subjects are informed of the denial and reason for such denial, as required, to meet the entity's objectives related to privacy.\par\sourcepage{61}
\end{aicpabox}
Selected AICPA point-of-focus headings (illustrative, not exhaustive): Responds to Data Controller Requests [P]; Authenticates Data Subjects' Identity [P][C]; Permits Data Subjects Access to Their Personal Information [P][C].

\begin{engineeringbox}
Authenticate applicable requests, locate personal information and respond within documented parameters.
\end{engineeringbox}
Suggested accountable function: Privacy Operations. Possible evidence: DSAR register; identity check; response record.

Illustrative test: Sample access requests for verification, completeness and response.

\Needspace{5\baselineskip}
\subsection*{P5.2 -- Correction and amendment}
\begin{aicpabox}
\textbf{P5.2}. The entity corrects, amends, or appends personal information based on information provided by data subjects and communicates such information to third parties, as committed or required, to meet the entity's objectives related to privacy. If a request for correction is denied, data subjects are informed of the denial and reason for such denial to meet the entity's objectives related to privacy.\par\sourcepage{62}
\end{aicpabox}
Selected AICPA point-of-focus headings (illustrative, not exhaustive): Responds to Data Controller Requests [P]; Communicates Denial of Access Requests [P][C]; Permits Data Subjects to Update or Correct Personal Information [P][C].

\begin{engineeringbox}
Support correction requests, document decisions and propagate changes as applicable.
\end{engineeringbox}
Suggested accountable function: Privacy Operations / Data. Possible evidence: Correction tickets; updated records; vendor notices.

Illustrative test: Trace completed and denied corrections with reason and propagation.

\section{Additional Criteria for Privacy: Disclosure and Notification}
\Needspace{5\baselineskip}
\subsection*{P6.1 -- Authorized disclosure}
\begin{aicpabox}
\textbf{P6.1}. The entity discloses personal information to third parties with the explicit consent of data subjects and such consent is obtained prior to disclosure to meet the entity's objectives related to privacy.\par\sourcepage{63}
\end{aicpabox}
Selected AICPA point-of-focus headings (illustrative, not exhaustive): Communicates Privacy Policies to Third Parties [P][C]; Discloses Personal Information Only When Appropriate [P][C]; Discloses Personal Information Only to Appropriate Third Parties [P][C].

\begin{engineeringbox}
Gate outbound personal information sharing by lawful authorization and documented purpose.
\end{engineeringbox}
Suggested accountable function: Privacy / Legal. Possible evidence: Data-sharing inventory; agreements; consent reference.

Illustrative test: Sample disclosures for authorized basis before transfer.

\Needspace{5\baselineskip}
\subsection*{P6.2 -- Disclosure accounting}
\begin{aicpabox}
\textbf{P6.2}. The entity creates and retains a complete, accurate, and timely record of authorized disclosures of personal information to meet the entity's objectives related to privacy.\par\sourcepage{63}
\end{aicpabox}
Selected AICPA point-of-focus headings (illustrative, not exhaustive): Creates and Retains Record of Authorized Disclosures [P][C].

\begin{engineeringbox}
Log approved personal information disclosures with recipient, data scope and date.
\end{engineeringbox}
Suggested accountable function: Privacy / Data Engineering. Possible evidence: Disclosure log; export records; reconciliation.

Illustrative test: Reconcile sampled data exports to complete authorization ledger.

\Needspace{5\baselineskip}
\subsection*{P6.3 -- Unauthorized disclosure recording}
\begin{aicpabox}
\textbf{P6.3}. The entity creates and retains a complete, accurate, and timely record of detected or reported unauthorized disclosures (including breaches) of personal information to meet the entity's objectives related to privacy.\par\sourcepage{64}
\end{aicpabox}
\begin{engineeringbox}
Record suspected and confirmed unauthorized disclosures with details and disposition.
\end{engineeringbox}
Suggested accountable function: Privacy / Incident Response. Possible evidence: Incident register; disclosure timeline; investigation records.

Illustrative test: Trace sampled events to documented records and closure rationale.

\Needspace{5\baselineskip}
\subsection*{P6.4 -- Privacy vendor commitments}
\begin{aicpabox}
\textbf{P6.4}. The entity obtains privacy commitments from vendors and other third parties who have access to personal information to meet the entity's objectives related to privacy. The entity assesses those parties' compliance on a periodic and as-needed basis and takes corrective action, if necessary.\par\sourcepage{64}
\end{aicpabox}
Selected AICPA point-of-focus headings (illustrative, not exhaustive): Evaluates Third-Party Compliance With Privacy Commitments [P][C]; Remediates Misuse of Personal Information by a Third Party [P][C]; Obtains Commitments to Report Unauthorized Disclosures [P][C].

\begin{engineeringbox}
Establish vendor privacy commitments, assess adherence and remediate misuse.
\end{engineeringbox}
Suggested accountable function: Privacy / Third-Party Risk. Possible evidence: DPA/contract clauses; vendor assessments; remediation log.

Illustrative test: Sample processors for commitment, review and corrective action.

\Needspace{5\baselineskip}
\subsection*{P6.5 -- Vendor incident notifications}
\begin{aicpabox}
\textbf{P6.5}. The entity obtains commitments from vendors and other third parties with access to personal information to notify the entity in the event of actual or suspected unauthorized disclosures of personal information. Such notifications are reported to appropriate personnel and acted on in accordance with established incident-response procedures to meet the entity's objectives related to privacy.\par\sourcepage{65}
\end{aicpabox}
Selected AICPA point-of-focus headings (illustrative, not exhaustive): Remediates Misuse of Personal Information by a Third Party [P][C]; Reports Actual or Suspected Unauthorized Disclosures [P][C].

\begin{engineeringbox}
Obtain escalation commitments and integrate third-party notifications into IR procedures.
\end{engineeringbox}
Suggested accountable function: Privacy / Third-Party Risk. Possible evidence: Notification terms; escalation contacts; vendor tests.

Illustrative test: Examine sampled vendor incidents and reporting timelines.

\Needspace{5\baselineskip}
\subsection*{P6.6 -- Privacy incident notification}
\begin{aicpabox}
\textbf{P6.6}. The entity provides notification of breaches and incidents to affected data subjects, regulators, and others to meet the entity's objectives related to privacy.\par\sourcepage{65}
\end{aicpabox}
Selected AICPA point-of-focus headings (illustrative, not exhaustive): Identifies Reporting Requirements [P][C]; Provides Notice of Breaches and Incidents [P][C].

\begin{engineeringbox}
Assess recipients and deadlines from applicable commitments and laws; document notification decisions.
\end{engineeringbox}
Suggested accountable function: Privacy / Legal. Possible evidence: Breach decision tree; counsel review; notices.

Illustrative test: Trace incidents to legal analysis and timely notification or rationale.

\Needspace{5\baselineskip}
\subsection*{P6.7 -- Accounting and disclosure to subject}
\begin{aicpabox}
\textbf{P6.7}. The entity provides data subjects with an accounting of the personal information held and disclosure of the data subjects' personal information, upon the data subjects' request, to meet the entity's objectives related to privacy.\par\sourcepage{65}
\end{aicpabox}
Selected AICPA point-of-focus headings (illustrative, not exhaustive): Responds to Data Controller Requests [P]; Identifies Types of Personal Information and Handling Process [P][C]; Captures, Identifies, and Communicates Requests for Information [P][C].

\begin{engineeringbox}
Provide appropriate accounting of held and disclosed personal information using verified records.
\end{engineeringbox}
Suggested accountable function: Privacy Operations. Possible evidence: DSAR package; inventory; disclosure accounting.

Illustrative test: Sample requests for completeness, identity verification and response.

\section{Additional Criteria for Privacy: Quality}
\Needspace{5\baselineskip}
\subsection*{P7.1 -- Personal data quality}
\begin{aicpabox}
\textbf{P7.1}. The entity collects and maintains accurate, up-to-date, complete, and relevant personal information to meet the entity's objectives related to privacy.\par\sourcepage{66}
\end{aicpabox}
Selected AICPA point-of-focus headings (illustrative, not exhaustive): Ensures Accuracy and Completeness of Personal Information [P][C]; Ensures Relevance of Personal Information [P][C].

\begin{engineeringbox}
Define validity, currency and completeness rules with correction and synchronization.
\end{engineeringbox}
Suggested accountable function: Data Governance. Possible evidence: Data-quality dashboard; reconciliations; corrected records.

Illustrative test: Sample key attributes for quality controls and exceptions.

\section{Additional Criteria for Privacy: Monitoring and Enforcement}
\Needspace{5\baselineskip}
\subsection*{P8.1 -- Privacy inquiry and dispute handling}
\begin{aicpabox}
\textbf{P8.1}. The entity implements a process for receiving, addressing, resolving, and communicating the resolution of inquiries, complaints, and disputes from data subjects and others and periodically monitors compliance to meet the entity's objectives related to privacy. Corrections and other necessary actions related to identified deficiencies are made or taken in a timely manner.\par\sourcepage{66}
\end{aicpabox}
Selected AICPA point-of-focus headings (illustrative, not exhaustive): Communicates to Data Subjects or Data Controllers [P][C]; Addresses Inquiries, Complaints, and Disputes [P][C]; Documents and Communicates Dispute Resolution and Recourse [P][C].

\begin{engineeringbox}
Offer intake, investigation, resolution, appeals and trend monitoring for privacy inquiries.
\end{engineeringbox}
Suggested accountable function: Privacy / Compliance. Possible evidence: Case register; SLA results; corrective actions.

Illustrative test: Select disputes and examine receipt, handling, decision and follow-up.

\chapter{Decision rules for application}
\begin{enumerate}
\item Start with the service organization system description and customer service commitments; explicitly define system boundaries and suppliers.
\item Identify which trust services categories are included, keeping the common criteria in scope.
\item For each applicable criterion, write the control objective and the associated risk scenario in organization-specific terms.
\item Identify the control owner, execution procedure, control frequency, and supporting technology; maintain the independence of monitoring where applicable.
\item Plan evidence at implementation time, including population-level exports, access control, retention, audit trail and reproducible timestamps.
\item Determine whether evidence shows design at a point in time or operation over a period; do not conflate the two.
\item For an identified criterion that management judges not applicable to the scoped system, document the facts, rationale and agreement with the practitioner.
\end{enumerate}

Cross-document keys: the identifier CC7.1, for example, is identical in the study guide, control mapping, readiness assessment, evidence matrix and testing handbook. The companion CSV provides one row for each of the 61 criteria.


\chapter*{Source and editorial protocol}
\addcontentsline{toc}{chapter}{Source and editorial protocol}
This independent document is grounded in the user-supplied \emph{AICPA TSP Section 100: 2017 Trust Services Criteria for Security, Availability, Processing Integrity, Confidentiality, and Privacy (With Revised Points of Focus---2022)}, 75-page PDF, especially paragraphs .01--.29, the criterion table on pp.~14--67, and the glossary. Criterion wording is quoted when labeled as such. Selected point-of-focus labels are aids to interpretation and may be incomplete. Engineering recommendations, role assignments, tools, tests, dates, scorecards, and example frequencies are author-developed examples and are not AICPA requirements. The 2022 revisions changed points of focus, not the 2017 criteria (source p.~3). The source states that an assessment of each point of focus is not required (p.~5). This document does not constitute a CPA attestation or legal advice.
\end{document}
